\documentclass[11pt,a4paper]{article}
\usepackage{amsmath,amssymb,amsthm}
\usepackage{mathtools}
\usepackage{bm}
\usepackage{booktabs}
\usepackage{array}
\usepackage{enumitem}
\usepackage{geometry}
\usepackage{xcolor}
\usepackage{tcolorbox}
\usepackage[pdfencoding=auto,psdextra]{hyperref}

\geometry{margin=2.5cm}

\newcommand{\E}{\mathbb{E}}
\newcommand{\R}{\mathbb{R}}
\newcommand{\N}{\mathcal{N}}
\newcommand{\KL}{\mathrm{KL}}
\newcommand{\TV}{\mathrm{TV}}
\newcommand{\norm}[1]{\left\lVert #1 \right\rVert}
\newcommand{\inner}[2]{\left\langle #1,\,#2 \right\rangle}
\newcommand{\sg}{\operatorname{sg}}
\newcommand{\dd}{\mathrm{d}}
\newcommand{\vold}{v_{\mathrm{old}}}
\newcommand{\vT}{v_T}
\newcommand{\vtheta}{v_\theta}
\DeclareMathOperator{\logit}{logit}

\theoremstyle{plain}
\newtheorem{theorem}{Theorem}[section]
\newtheorem{proposition}[theorem]{Proposition}
\newtheorem{lemma}[theorem]{Lemma}
\newtheorem{corollary}[theorem]{Corollary}

\theoremstyle{definition}
\newtheorem{definition}[theorem]{Definition}
\newtheorem{remark}[theorem]{Remark}

\title{\textbf{Five Formalizations of Proximal On-Policy Distillation for Flow Models}\\[0.3em]
\large From deterministic field matching to stochastic transition and marginal mixtures}
\author{Jayce-Ping}
\date{2026-07-30}

\begin{document}
\maketitle

\begin{abstract}
Flow models admit both deterministic ODE sampling and stochastic SDE sampling.  A proximal
distillation rule built from Gaussian transition likelihoods is therefore exact only for the SDE
case; its deterministic limit is singular.  This note develops five concrete arms under one
flow-matching notation.  Arm A0 is direct ODE teacher-field regression.  Arm A1 adds an explicit
field-space trust region with a local one-step transport interpretation.  Arm A2 constructs a
teacher gate from forward flow-matching prediction risk; it is solver-agnostic but is a predictive
surrogate rather than a state-density ratio.  Arm A3 is the existing SDE transition-mixture P-OPD,
whose detached-responsibility loss matches the exact mixture-KL gradient locally at the behavior
student.  Arm A4 uses old and teacher trajectories in a two-source conditional flow-matching
objective and exactly realizes their mixture of time marginals in the population limit.

The derivations stay at the level needed for flow-matching practice: every arm specifies its
sampling distribution, target, loss, gradient, scale convention, exactness claim, and failure
boundary.  In particular, plain ODE MSE is not an unscaled zero-noise KL; it is a quadratic
field/transport geometry, or equivalently a rescaled small-noise control-energy limit.
\end{abstract}

\tableofcontents
\newpage

%% ============================================================
\section{Common notation and reusable results}
\label{sec:common}
%% ============================================================

\subsection{Flow, rollout, and scale conventions}

We parameterize time so that sampling integrates from $t=0$ (base noise) to $t=1$ (generated
sample).  This is only a change of orientation from schedulers that denoise from $1$ to $0$.  Let
$C$ denote conditioning information and $Z\sim\rho_{\mathrm{base}}(\cdot\mid C)$ the initial
noise.  At outer iteration $k$, the frozen behavior student has vector field
$\vold=v_{\theta_k}$ and rollout
\begin{equation}
\frac{\dd X_t^{\mathrm{old}}}{\dd t}
  =\vold(t,X_t^{\mathrm{old}},C),
\qquad
X_0^{\mathrm{old}}=Z .
\label{eq:old-ode}
\end{equation}
The frozen teacher field is $\vT$, and $\vtheta$ is the trainable field.  Their flow maps and
time marginals are denoted by $\Phi_t^{\mathrm{old}},\Phi_t^T,\Phi_t^\theta$ and
$\rho_{\mathrm{old},t},\rho_{T,t},\rho_{\theta,t}$.

Let $\omega(t)>0$ be a normalized training-time density,
$\int_0^1\omega(t)\dd t=1$.  The old-occupancy expectation is
\begin{equation}
\E_{\mathrm{old}}[f]
  :=
  \E_{C,Z}\int_0^1
  \omega(t)
  f(t,X_t^{\mathrm{old}},C)\dd t .
\label{eq:old-occupancy}
\end{equation}
The state $X_t^{\mathrm{old}}$ is detached during optimization.  Thus
\eqref{eq:old-occupancy} describes fixed-input field regression; it does not differentiate
through the ODE solver.

For a detached positive-definite metric $G_t$, define
\begin{equation}
\norm{a}_{G_t}^2=a^\top G_ta,
\qquad
\norm{a}_{\mathrm{old},G}^2
  =\E_{\mathrm{old}}\!\left[\norm{a(t,X_t^{\mathrm{old}},C)}_{G_t}^2\right].
\label{eq:field-norm}
\end{equation}
The default practical choice is $G_t=I$.  A covariance-derived metric
$G_t=A_t^{-1}$ is meaningful only when a reference covariance $A_t$ has been explicitly declared;
an ODE does not supply one automatically.

If the latent event has dimension $D$, then
\begin{equation}
\operatorname{RMS}_{G_t}(a)^2
  :=\frac{1}{D}\norm{a}_{G_t}^2,
\qquad
\norm{a}_{G_t}^2
  =D\,\operatorname{RMS}_{G_t}(a)^2.
\label{eq:rms-vs-sum}
\end{equation}
The event sum is the scale used by a joint Gaussian density.  The RMS or event mean is easier to
compare across resolutions.  A trust radius or temperature must state which convention it uses;
the two differ by $D$ in squared scale.

\subsection{Three different objects}

The five arms act on three related but non-interchangeable objects:
\begin{center}
\begin{tabular}{@{}lll@{}}
\toprule
object & notation & what is compared\\
\midrule
one-step conditional kernel & $K(\dd y\mid x,t)$ & stochastic next-state law\\
vector field / flow map & $v(t,x)$, $\Phi_t(z)$ & deterministic transport rule\\
time marginal & $\rho_t=(\Phi_t)_\#\rho_{\mathrm{base}}$ & distribution of states at time $t$\\
\bottomrule
\end{tabular}
\end{center}
An SDE transition KL belongs to the first row.  ODE velocity MSE belongs to the second.  A
marginal-mixture construction belongs to the third.  Replacing one row by another can produce a
useful surrogate, but it is not an algebraic identity.

\subsection{Why deterministic transition KL is singular}

For one numerical step of size $h$, let
\begin{equation}
F_j^h(x)=x+h\,v_j(t,x),
\qquad
K_j(\dd y\mid x)=\delta_{F_j^h(x)}(\dd y),
\qquad
j\in\{\mathrm{old},T\}.
\label{eq:deterministic-kernels}
\end{equation}

\begin{proposition}[No smooth conditional-KL trust region for unequal ODE maps]
\label{prop:dirac-kl}
Assume $F_{\mathrm{old}}^h(x)\neq F_T^h(x)$ and define
\[
Q_\alpha=(1-\alpha)K_{\mathrm{old}}+\alpha K_T,
\qquad \alpha\in(0,1).
\]
Then
\begin{equation}
\KL(\delta_y\Vert Q_\alpha)
=
\begin{cases}
-\log(1-\alpha), & y=F_{\mathrm{old}}^h(x),\\
-\log\alpha, & y=F_T^h(x),\\
+\infty, & \text{otherwise}.
\end{cases}
\label{eq:dirac-mixture-kl}
\end{equation}
Therefore the objective is finite at the old map but infinite after a generic arbitrarily small
output perturbation.  It has no local teacher-pull gradient.
\end{proposition}

\begin{proof}
$Q_\alpha$ places mass only at the two displayed atoms.  The Radon--Nikodym derivative of
$\delta_y$ with respect to $Q_\alpha$ is $1/(1-\alpha)$ at the old atom, $1/\alpha$ at the
teacher atom, and does not exist elsewhere.  Substitution into the definition of KL gives
\eqref{eq:dirac-mixture-kl}.
\end{proof}

JS divergence is finite here, but it is piecewise constant as the atoms move and supplies no
useful local teacher-pull gradient.  A deterministic PPO-style likelihood ratio is likewise
undefined.  The problem is the singular support, not merely the choice of divergence.

\subsection{What does survive the small-noise limit}

Smooth two outputs with a shared covariance $\varepsilon A$, where $A\succ0$:
\begin{equation}
p_a^\varepsilon=\N(a,\varepsilon A),
\qquad
p_b^\varepsilon=\N(b,\varepsilon A).
\end{equation}
The equal-covariance Gaussian identity gives
\begin{equation}
\KL(p_a^\varepsilon\Vert p_b^\varepsilon)
  =\frac{1}{2\varepsilon}\norm{a-b}_{A^{-1}}^2,
\qquad
2\varepsilon\,
\KL(p_a^\varepsilon\Vert p_b^\varepsilon)
  =\norm{a-b}_{A^{-1}}^2 .
\label{eq:rescaled-gaussian-kl}
\end{equation}
The unscaled KL diverges as $\varepsilon\to0$ for $a\neq b$, but the rescaled KL is exactly a
quadratic transport cost.

For Euler means $\mu_j=x+h\,v_j$ and diffusion covariance
$\Sigma=\varepsilon h A$,
\begin{equation}
\KL\!\left(
  \N(\mu_1,\varepsilon hA)
  \,\middle\Vert\,
  \N(\mu_2,\varepsilon hA)
\right)
=\frac{h}{2\varepsilon}\norm{v_1-v_2}_{A^{-1}}^2 .
\label{eq:euler-kl-control}
\end{equation}
Let $P_1^{\mathrm{grid}}$ and $P_2^{\mathrm{grid}}$ be the two Markov trajectory laws on the same
grid and with the same initial law.  The KL chain rule and \eqref{eq:euler-kl-control} give
\begin{align}
\varepsilon\KL(P_1^{\mathrm{grid}}\Vert P_2^{\mathrm{grid}})
&=
\frac12
\E_{P_1^{\mathrm{grid}}}
\sum_\ell h_\ell
\norm{
  v_1(t_\ell,X_\ell)-v_2(t_\ell,X_\ell)
}_{A_\ell^{-1}}^2
\nonumber\\
&\longrightarrow
\frac12
\E_{P_1}
\int_0^1
\norm{v_1(t,X_t)-v_2(t,X_t)}_{A_t^{-1}}^2\dd t .
\label{eq:control-energy-limit}
\end{align}
Equation~\eqref{eq:control-energy-limit} is the principled bridge from SDE KL to ODE MSE.  It
preserves a rescaled pairwise control cost; it does \emph{not} preserve the transition-mixture
responsibility gate, as Section~\ref{sec:a3} will show.

\subsection{A field-space trust region controls ODE displacement}

\begin{proposition}[Old-occupancy field bound]
\label{prop:gronwall}
Assume $\vtheta(t,\cdot,c)$ is $L$-Lipschitz for every $(t,c)$.  Couple the candidate and old ODEs
with the same $(C,Z)$.  Then
\begin{align}
\norm{X_1^\theta-X_1^{\mathrm{old}}}
&\le
\int_0^1 e^{L(1-s)}
\norm{
  \vtheta(s,X_s^{\mathrm{old}},C)
  -\vold(s,X_s^{\mathrm{old}},C)
}\dd s,
\label{eq:pathwise-gronwall}\\
\E\norm{X_1^\theta-X_1^{\mathrm{old}}}^2
&\le
e^{2L}
\int_0^1
\E\norm{
  \vtheta(s,X_s^{\mathrm{old}},C)
  -\vold(s,X_s^{\mathrm{old}},C)
}^2\dd s .
\label{eq:mean-gronwall}
\end{align}
Consequently,
\begin{equation}
W_2^2(\rho_{\theta,1},\rho_{\mathrm{old},1})
\le
\E\norm{X_1^\theta-X_1^{\mathrm{old}}}^2 .
\label{eq:w2-coupling}
\end{equation}
If additionally
\begin{equation}
\omega(t)\ge\omega_*>0,
\qquad
G_t\succeq g_*I
\label{eq:field-bound-assumptions}
\end{equation}
for all trained times, then the field norm in \eqref{eq:field-norm} controls the endpoint:
\begin{equation}
W_2^2(\rho_{\theta,1},\rho_{\mathrm{old},1})
\le
\frac{e^{2L}}{\omega_*g_*}
\norm{\vtheta-\vold}_{\mathrm{old},G}^2 .
\label{eq:field-norm-endpoint}
\end{equation}
\end{proposition}

\begin{proof}
Add and subtract
$\vtheta(t,X_t^{\mathrm{old}},C)$ in the difference of the two ODEs.  The candidate-field
Lipschitz bound and Gr\"onwall's inequality give \eqref{eq:pathwise-gronwall}.  Cauchy--Schwarz
gives \eqref{eq:mean-gronwall}.  The common-noise pairing is a valid coupling of the two endpoint
laws, proving \eqref{eq:w2-coupling}.  Finally,
\emph{under the additional lower bounds} in \eqref{eq:field-bound-assumptions},
\[
\norm{\vtheta-\vold}_{\mathrm{old},G}^2
\ge
\omega_*g_*
\int_0^1
\E\norm{\vtheta(t,X_t^{\mathrm{old}})-\vold(t,X_t^{\mathrm{old}})}^2\dd t,
\]
which gives \eqref{eq:field-norm-endpoint}.
\end{proof}

\begin{remark}
The bound controls departure from the \emph{old} rollout because the field error is evaluated at
old states.  Replacing it by a candidate--teacher endpoint guarantee requires candidate-occupancy
teacher error, or an additional coverage/Lipschitz argument.  Old-state MSE alone does not give a
global teacher-distribution convergence theorem.
\end{remark}

%% ============================================================
\section{Arm A0: Direct ODE-OPD}
\label{sec:a0}
%% ============================================================

\subsection{Objective and gradient}

Direct ODE-OPD queries the teacher on states visited by the behavior student and regresses the
current field to the teacher:
\begin{equation}
\mathcal L_{\mathrm{A0}}(\theta)
  =
  \frac12
  \norm{\vtheta-\vT}_{\mathrm{old},G}^2
  =
  \frac12
  \E_{\mathrm{old}}\!\left[
    \norm{\vtheta(t,X_t^{\mathrm{old}},C)
          -\vT(t,X_t^{\mathrm{old}},C)}_{G_t}^2
  \right].
\label{eq:a0-loss}
\end{equation}
This is the same direct-matching family as DiffusionOPD/XOPD~\cite{diffusionopd}: rollout is
on-policy, while optimization is closed-form regression and does not require transition
likelihoods.  Their deterministic transition-mean objective corresponds specifically to the
$h_\ell^2$-weighted velocity loss derived below, rather than to an arbitrary choice of
$(\omega,G)$ in \eqref{eq:a0-loss}.

\begin{proposition}[A0 is old-occupancy weighted field regression]
\label{prop:a0-projection}
In unrestricted output space, every attained population minimizer of \eqref{eq:a0-loss} satisfies
\begin{equation}
v^\star(t,x,c)=\vT(t,x,c)
\label{eq:a0-unrestricted}
\end{equation}
almost everywhere under the old occupation measure.  In a model class $\mathcal V$, any attained
global solution is the weighted least-squares fit
\begin{equation}
v_{\mathcal V}^\star
\in
\arg\min_{v\in\mathcal V}
\norm{v-\vT}_{\mathrm{old},G}^2 .
\label{eq:a0-finite-capacity}
\end{equation}
Its parameter gradient is
\begin{equation}
\nabla_\theta\mathcal L_{\mathrm{A0}}
=
\E_{\mathrm{old}}\!\left[
  J_\theta^\top G_t(\vtheta-\vT)
\right],
\qquad
J_\theta=\partial_\theta\vtheta .
\label{eq:a0-gradient}
\end{equation}
\end{proposition}

\begin{proof}
Positive definiteness gives
$\norm{v-\vT}_{G_t}^2\ge0$, with equality exactly at $v=\vT$.  Integration preserves this
characterization almost everywhere under the occupation measure.  Restricting the minimization
to $\mathcal V$ gives \eqref{eq:a0-finite-capacity}; differentiating the detached-input quadratic
gives \eqref{eq:a0-gradient}.
\end{proof}

For a linear output model $v_\theta=B(t,x,c)\theta$, the finite-capacity optimum obeys the weighted
normal equations
\begin{equation}
\E_{\mathrm{old}}[B^\top G_tB]\,\theta^\star
=
\E_{\mathrm{old}}[B^\top G_t\vT].
\label{eq:a0-normal-equations}
\end{equation}
For a nonlinear network, stationarity only implies tangent-space orthogonality
\begin{equation}
\E_{\mathrm{old}}\!\left[
  J_{\theta^\star}^{\top}G_t(v_{\theta^\star}-\vT)
\right]=0;
\label{eq:a0-tangent}
\end{equation}
it does not imply existence or uniqueness of a global least-squares solution.

\subsection{The practical loss space}

For an Euler step and a linear flow-matching path,
\begin{equation}
\mu_\theta-\mu_T
  =h(\vtheta-\vT),
\qquad
\widehat x_0^\theta-\widehat x_0^T
  =s(\vtheta-\vT).
\label{eq:a0-spaces}
\end{equation}
The plus sign follows from the common noise-to-data time orientation:
$x_s=x_0-sv$ and hence $\widehat x_0=x_s+sv$.  A native data-to-noise predictor
$u=\varepsilon-x_0=-v$ uses the opposite sign; the squared-loss scale below is unchanged.
Therefore velocity, next-state mean, and clean-prediction MSE differ only by timestep weights:
\begin{equation}
\mathcal L_v:\mathcal L_{x_t}:\mathcal L_{x_0}
\quad\Longleftrightarrow\quad
1:h^2:s^2 .
\label{eq:a0-space-weights}
\end{equation}
At one fixed timestep these factors only rescale the gradient.  Across a trajectory they define
different quadrature rules.  In particular,
\begin{equation}
\int_0^1\norm{\Delta v_t}^2\dd t
\approx
\sum_\ell h_\ell\norm{\Delta v_\ell}^2
=
\sum_\ell \frac{\norm{\Delta\mu_\ell}^2}{h_\ell},
\label{eq:a0-control-quadrature}
\end{equation}
whereas an unnormalized sum of next-state MSE uses $h_\ell^2$ and tends to zero under grid
refinement.

\subsection{What A0 does and does not guarantee}

If the fitted field remains close to the old field, Proposition~\ref{prop:gronwall} turns measured
field displacement into a bound on rollout displacement.  A0 itself, however, contains no
old-field penalty.  The triangle inequality only gives
\begin{equation}
\norm{\vtheta-\vold}_{\mathrm{old},G}
\le
\norm{\vtheta-\vT}_{\mathrm{old},G}
+
\norm{\vT-\vold}_{\mathrm{old},G},
\label{eq:a0-no-trust}
\end{equation}
which can be large when the teacher is far from the student.

The loss also constrains only old-occupancy regions.  Low-density or teacher-only states receive
little or no supervision in a finite batch.  Refreshing the student rollout after every outer
iteration improves coverage adaptively, but does not turn a local regression result into a global
distribution contraction.

\begin{tcolorbox}[colback=gray!5,colframe=gray!55,title=Exact interpretation of Arm A0]
A0 is on-policy \emph{field regression}: the inputs come from the current behavior flow.  It is
not proximal merely because those inputs are on-policy, and it is not an ODE transition KL because
unequal deterministic kernels are singular by Proposition~\ref{prop:dirac-kl}.
\end{tcolorbox}

%% ============================================================
\section{Arm A1: Transport-Projected ODE-OPD}
\label{sec:a1}
%% ============================================================

\subsection{Quadratic proximal target}

Arm A1 explicitly anchors the field to the behavior student:
\begin{equation}
\mathcal L_{\mathrm{A1},\alpha}(v)
=
\frac{\alpha}{2}\norm{v-\vT}_{\mathrm{old},G}^2
+
\frac{1-\alpha}{2}\norm{v-\vold}_{\mathrm{old},G}^2,
\qquad
\alpha\in[0,1].
\label{eq:a1-loss}
\end{equation}

\begin{proposition}[Barycentric target]
\label{prop:a1-barycenter}
Define
\begin{equation}
v_\alpha=(1-\alpha)\vold+\alpha\vT,
\qquad
D_k=\norm{\vT-\vold}_{\mathrm{old},G}.
\label{eq:a1-target}
\end{equation}
Then
\begin{equation}
\mathcal L_{\mathrm{A1},\alpha}(v)
=
\frac12\norm{v-v_\alpha}_{\mathrm{old},G}^2
+
\frac{\alpha(1-\alpha)}{2}D_k^2 .
\label{eq:a1-square-completion}
\end{equation}
Thus $v_\alpha$ is the unrestricted optimum.  In a finite neural class, training is a weighted
least-squares fit to $v_\alpha$, not a parameter interpolation between the two networks.
\end{proposition}

\begin{proof}
Apply pointwise
\[
\alpha\norm{v-a}_G^2+(1-\alpha)\norm{v-b}_G^2
=
\norm{v-\alpha a-(1-\alpha)b}_G^2
+
\alpha(1-\alpha)\norm{a-b}_G^2
\]
and average under the old occupancy.
\end{proof}

\subsection{Equivalent trust-budget projection}

Consider the explicit field-space trust problem
\begin{equation}
\min_v\ \frac12\norm{v-\vT}_{\mathrm{old},G}^2
\quad\text{subject to}\quad
\frac12\norm{v-\vold}_{\mathrm{old},G}^2\le\delta .
\label{eq:a1-constrained}
\end{equation}

\begin{proposition}[Closed-form trust projection]
\label{prop:a1-trust-projection}
Let $R=\sqrt{2\delta}$.  If $D_k>0$, the unrestricted solution of
\eqref{eq:a1-constrained} is
\begin{equation}
v_R
=
\vold+\alpha_R(\vT-\vold),
\qquad
\alpha_R=\min\!\left\{1,\frac{R}{D_k}\right\}.
\label{eq:a1-alpha-radius}
\end{equation}
If $D_k=0$, $v_R=\vold=\vT$.  Conversely, a fixed $\alpha$ corresponds to
\begin{equation}
R_\alpha=\alpha D_k,
\qquad
\delta_\alpha=\frac{\alpha^2D_k^2}{2}.
\label{eq:a1-relative-radius}
\end{equation}
\end{proposition}

\begin{proof}
Write $u=v-\vold$ and $d=\vT-\vold$.  The objective is
\[
\norm{u-d}_{\mathrm{old},G}^2
=\norm{u}_{\mathrm{old},G}^2+D_k^2
-2\inner{u}{d}_{\mathrm{old},G}.
\]
For fixed $\norm{u}_{\mathrm{old},G}$, Cauchy--Schwarz makes $u$ parallel to $d$ optimal.
The best allowed length is $\min\{R,D_k\}$, giving \eqref{eq:a1-alpha-radius}.
\end{proof}

A fixed $\alpha$ is therefore a \emph{relative} trust radius: the absolute update grows with the
teacher--old gap.  If a fixed absolute field RMS radius $r$ is desired, use
\begin{align}
d_k^{\mathrm{RMS}}
&=
\sqrt{\E_{\mathrm{old}}[
  \operatorname{RMS}_{G_t}(\vT-\vold)^2
]},
\nonumber\\
\alpha_k
&=
\begin{cases}
1, & d_k^{\mathrm{RMS}}=0,\\
\min\{1,r/d_k^{\mathrm{RMS}}\}, & d_k^{\mathrm{RMS}}>0.
\end{cases}
\label{eq:a1-rms-alpha}
\end{align}

\subsection{Wasserstein projection of a stochastic branch target}

At one old state, define the old and teacher Euler outputs
\begin{equation}
y_{\mathrm{old}}=x+h\vold(t,x),
\qquad
y_T=x+h\vT(t,x),
\end{equation}
and the branch mixture
\begin{equation}
Q_{\alpha,x}
=(1-\alpha)\delta_{y_{\mathrm{old}}}
+\alpha\delta_{y_T}.
\label{eq:a1-branch-mixture}
\end{equation}

\begin{proposition}[Best deterministic quadratic-transport projection]
\label{prop:a1-wasserstein}
For the quadratic ground metric $G_t$,
\begin{align}
\arg\min_y W_{2,G_t}^2(\delta_y,Q_{\alpha,x})
&=(1-\alpha)y_{\mathrm{old}}+\alpha y_T,
\label{eq:a1-wasserstein-minimizer}\\
\min_y W_{2,G_t}^2(\delta_y,Q_{\alpha,x})
&=\alpha(1-\alpha)
  \norm{y_T-y_{\mathrm{old}}}_{G_t}^2 .
\label{eq:a1-irreducible}
\end{align}
\end{proposition}

\begin{proof}
The only coupling from $\delta_y$ to $Q_{\alpha,x}$ sends masses $1-\alpha$ and $\alpha$ to the
two atoms, so
\[
W_{2,G_t}^2(\delta_y,Q_{\alpha,x})
=(1-\alpha)\norm{y-y_{\mathrm{old}}}_{G_t}^2
+\alpha\norm{y-y_T}_{G_t}^2.
\]
Completing the square yields \eqref{eq:a1-wasserstein-minimizer} and
\eqref{eq:a1-irreducible}.
\end{proof}

For Euler maps, the minimizer is $x+h v_\alpha$ and the irreducible cost is
\begin{equation}
\alpha(1-\alpha)h^2\norm{\vT-\vold}_{G_t}^2 .
\label{eq:a1-euler-cost}
\end{equation}
Thus A1 is the best deterministic approximation to a local stochastic branch mixture under the
chosen quadratic metric.  It does \emph{not} reproduce the two-atom law.

\subsection{Why one update is only a scaled A0 update}

\begin{proposition}[First-step anchor degeneracy]
\label{prop:a1-first-step}
Assume exact replay, so $v_{\theta_k}=\vold$ on every stored state.  Then
\begin{equation}
\left.
\nabla_\theta\mathcal L_{\mathrm{A1},\alpha}
\right|_{\theta_k}
=
\alpha
\left.
\nabla_\theta\mathcal L_{\mathrm{A0}}
\right|_{\theta_k}.
\label{eq:a1-first-gradient}
\end{equation}
\end{proposition}

\begin{proof}
The old-anchor gradient contains
$J_{\theta_k}^{\top}G_t(v_{\theta_k}-\vold)$, which is pointwise zero.  Only the
teacher term remains, multiplied by $\alpha$.
\end{proof}

For plain SGD, one step is therefore A0 with learning rate $\alpha\eta$.  Adam moments, clipping,
weight decay, and the optimizer epsilon can prevent the final parameter update from scaling
exactly by $\alpha$, but the loss gradient identity remains exact.  If output-space gradient
descent with $G_t=I$ and step size $0<\eta<2$ is repeated against a frozen old target, then
\begin{equation}
v^{(m)}
=
\vold
+\alpha\bigl[1-(1-\eta)^m\bigr](\vT-\vold).
\label{eq:a1-inner-steps}
\end{equation}
The old anchor becomes active after the first update and the limit is $v_\alpha$.  Genuine
first-update trust enforcement instead requires explicit projection, a constrained line search,
or post-update stopping based on the measured field displacement.

\begin{tcolorbox}[colback=gray!5,colframe=gray!55,title=Practical meaning of Arm A1]
A1 supplies a real geometric trust region only when the radius is enforced or the anchored
subproblem is optimized beyond its first gradient.  With one on-policy optimizer step, a fixed
$\alpha$ is only a teacher-gradient scale.  Flow-DPPO~\cite{flowdppo} motivates controlling
current--old divergence; A1 is a field-space analogue with a local one-step transport
interpretation, not a deterministic likelihood ratio or a full-flow-map projection.
\end{tcolorbox}

%% ============================================================
\section{Arm A2: Forward-Risk P-OPD}
\label{sec:a2}
%% ============================================================

\subsection{Forward-process evidence}

Arm A2 keeps efficient ODE rollout but constructs a supervised flow-matching test after the
rollout.  Let $x_{\mathrm{data}}=X_1^{\mathrm{old}}$ be the old ODE endpoint, draw
$s\sim\omega_{\mathrm{FM}}$ on $[0,1]$ and
$\varepsilon\sim\N(0,I_D)$, and form
\begin{equation}
x_s=(1-s)x_{\mathrm{data}}+s\varepsilon
    =x_{\mathrm{data}}+s u,
\qquad
u=\varepsilon-x_{\mathrm{data}}.
\label{eq:a2-forward-path}
\end{equation}
Let $z=(x_s,s,C)$ and evaluate frozen predictions
$f_{\mathrm{old}}(z)$ and $f_T(z)$ of the native data-to-noise target $u$ on the same input.
Every A2 expectation below is over $(C,Z,s,\varepsilon)$ under this construction.  This is the
solver-decoupled forward-flow-matching construction~\cite{flowmatching} shared by DPO,
DiffusionNFT, AWM, DGPO, and CRD.

To make event scaling explicit, define
\begin{equation}
E_j(z,u)
=
\frac{1}{2c_D}\norm{u-f_j(z)}^2,
\qquad
c_D=
\begin{cases}
1, & \text{event sum},\\
D, & \text{event mean},
\end{cases}
\label{eq:a2-energy}
\end{equation}
and
\begin{equation}
A_T=E_{\mathrm{old}}-E_T,
\qquad
\gamma_T
=
\sigma\!\left(
  \logit\alpha+\frac{A_T}{\tau_s}
\right).
\label{eq:a2-gate}
\end{equation}
Positive $A_T$ means that the teacher predicts the sampled forward target better than the old
student.

\subsection{Exact auxiliary Gaussian interpretation}

\begin{proposition}[When the forward-risk gate is a Bayes responsibility]
\label{prop:a2-gaussian}
Fix $z$ and introduce a component label $B\in\{\mathrm{old},T\}$ with
$\Pr(B=T)=\alpha$.  Declare the predictive likelihood
\begin{equation}
p_j(u\mid z)
=
(2\pi\sigma_r^2(s))^{-D/2}
\exp\!\left[
  -\frac{\norm{u-f_j(z)}^2}{2\sigma_r^2(s)}
\right].
\label{eq:a2-predictive-model}
\end{equation}
Then
\begin{equation}
\Pr(B=T\mid u,z)
=
\sigma\!\left(
  \logit\alpha
  +\frac{c_D A_T}{\sigma_r^2(s)}
\right).
\label{eq:a2-posterior}
\end{equation}
Therefore \eqref{eq:a2-gate} is exactly this posterior when
\begin{equation}
\tau_s=\frac{\sigma_r^2(s)}{c_D}.
\label{eq:a2-exact-temperature}
\end{equation}
\end{proposition}

\begin{proof}
The common Gaussian normalizers cancel and
\[
\log\frac{p_T(u\mid z)}{p_{\mathrm{old}}(u\mid z)}
=
\frac{
  \norm{u-f_{\mathrm{old}}}^2-\norm{u-f_T}^2
}{2\sigma_r^2(s)}
=
\frac{c_DA_T}{\sigma_r^2(s)}.
\]
Bayes' rule for the two-component mixture gives \eqref{eq:a2-posterior}.
\end{proof}

The covariance $\sigma_r^2(s)I$ in Proposition~\ref{prop:a2-gaussian} is a \emph{declared
predictive-residual model}; it is not automatically equal to the variance used to create $x_s$.
Unequal component covariances would add log-determinant and quadratic terms that a plain
square-loss difference does not contain.  Using an event mean with an unchanged numerical
temperature is equivalent to inflating the predictive covariance by $D$.

\subsection{What risk difference estimates}

\begin{proposition}[Conditional square-loss decomposition]
\label{prop:a2-risk-decomposition}
Let
\begin{equation}
m_{\mathrm{old}}(z)=\E[u\mid z],
\qquad
V_{\mathrm{old}}(z)
=
\E\!\left[
  \norm{u-m_{\mathrm{old}}(z)}^2\mid z
\right],
\label{eq:a2-conditional-mean}
\end{equation}
where the expectation is under old-generated endpoints and fresh forward noise.  Then
\begin{equation}
\E[E_j\mid z]
=
\frac{1}{2c_D}
\left(
  V_{\mathrm{old}}(z)
  +\norm{f_j(z)-m_{\mathrm{old}}(z)}^2
\right)
\label{eq:a2-bias-variance}
\end{equation}
and
\begin{equation}
\E[A_T\mid z]
=
\frac{
  \norm{f_{\mathrm{old}}-m_{\mathrm{old}}}^2
  -\norm{f_T-m_{\mathrm{old}}}^2
}{2c_D}.
\label{eq:a2-expected-advantage}
\end{equation}
\end{proposition}

\begin{proof}
Write
$u-f_j=(u-m_{\mathrm{old}})+(m_{\mathrm{old}}-f_j)$.  The conditional cross term vanishes because
$\E[u-m_{\mathrm{old}}\mid z]=0$.  Subtracting the two decompositions cancels the same
irreducible conditional variance.
\end{proof}

Thus $A_T$ is a one-sample unbiased estimate of the conditional \emph{risk difference}.  The
nonlinear gate is not an unbiased estimate of the sigmoid of the mean risk difference:
\begin{equation}
\E[\gamma_T\mid z]
\neq
\sigma\!\left(
  \logit\alpha+\frac{\E[A_T\mid z]}{\tau_s}
\right)
\quad\text{in general}.
\label{eq:a2-sigmoid-bias}
\end{equation}
Also, the old ODE field need not equal $m_{\mathrm{old}}(z)$ after an endpoint has been re-noised
with fresh $\varepsilon$.  The construction compares two predictors on old-generated data; it
does not assume that the old predictor is Bayes-optimal for the new coupling.

\subsection{Gated projected target and trust radius}

Let $r_s>0$ be an event-$L^2$ output radius in metric $G_s$.  A per-dimension RMS radius
$r_s^{\mathrm{RMS}}$ corresponds to $r_s=\sqrt D\,r_s^{\mathrm{RMS}}$.  Define
\begin{equation}
d(z)=f_T(z)-f_{\mathrm{old}}(z),
\qquad
c(z)=
\begin{cases}
1, & \norm{d(z)}_{G_s}=0,\\
\min\{1,r_s/\norm{d(z)}_{G_s}\}, & \norm{d(z)}_{G_s}>0.
\end{cases}
\label{eq:a2-projection-coef}
\end{equation}
The detached target and loss are
\begin{align}
f_\gamma(z,u)
&=
f_{\mathrm{old}}(z)
+\sg(\gamma_T)c(z)d(z),
\label{eq:a2-target}\\
\mathcal L_{\mathrm{A2}}(\theta)
&=
\E\!\left[
  \frac{1}{2c_D}
  \norm{f_\theta(z)-f_\gamma(z,u)}_{G_s}^2
\right].
\label{eq:a2-loss}
\end{align}
Since $0\le\gamma_T,c\le1$,
\begin{equation}
\norm{f_\gamma-f_{\mathrm{old}}}_{G_s}\le r_s.
\label{eq:a2-target-trust}
\end{equation}
At exact replay, $f_{\theta_{\mathrm{old}}}=f_{\mathrm{old}}$, and
\begin{equation}
\left.
\nabla_\theta\mathcal L_{\mathrm{A2}}
\right|_{\theta_{\mathrm{old}}}
=
-\frac{1}{c_D}
\E\!\left[
  \gamma_Tc(z)
  J_{\mathrm{old}}^\top G_s
  (f_T-f_{\mathrm{old}})
\right].
\label{eq:a2-local-gradient}
\end{equation}
The gate decides whether the teacher is useful on this supervised forward target; $c(z)$ decides
how far the target is allowed to move.

For any $\delta\in(0,\tfrac12)$,
\begin{align}
\gamma_T\le\delta
&\iff
A_T\le
\tau_s\bigl(\logit\delta-\logit\alpha\bigr),\\
\gamma_T\ge1-\delta
&\iff
A_T\ge
\tau_s\bigl(\logit(1-\delta)-\logit\alpha\bigr).
\label{eq:a2-saturation}
\end{align}
An event sum normally scales as $D$ at a fixed per-coordinate risk gap, so a fixed $\tau_s$
becomes increasingly binary with resolution.  An event mean is stable across $D$ but changes the
auxiliary Gaussian model.  A practical implementation should log both $A_T$ and $A_T/D$, gate
entropy and saturation, $\norm{d}_{\mathrm{RMS}}$, $\norm{d}_{2}$, $c(z)$, and the realized
post-update field displacement.

\begin{tcolorbox}[colback=yellow!5,colframe=yellow!45!black,title=Exact boundary of Arm A2]
The Gaussian statement concerns an auxiliary likelihood for the target $u$ conditioned on $z$.
Neither $\rho_T(x_s)$ nor $\rho_{\mathrm{old}}(x_s)$ appears.  Therefore the gate is predictive
risk evidence on old-generated samples, not a teacher/old state-density ratio and not an exact
marginal-mixture objective.  Timestep centering, running-scale normalization, or a calibrated
temperature can be useful, but they define an explicit surrogate.
\end{tcolorbox}

%% ============================================================
\section{Arm A3: SDE Transition-Mixture P-OPD}
\label{sec:a3}
%% ============================================================

\subsection{Exact one-step mixture}

Unlike Sections~\ref{sec:a0}--\ref{sec:a2}, Arm A3 uses a frozen behavior \emph{SDE} rollout.
On a grid, sample
\begin{equation}
X_{\ell+1}^{\mathrm{old}}
\sim
p_{\mathrm{old},\ell}(
  \,\cdot\mid X_\ell^{\mathrm{old}},C
),
\qquad
J\sim\nu_{\mathrm{train}},
\label{eq:a3-behavior-rollout}
\end{equation}
and let $\E_{\mathrm{SDE-old}}$ denote expectation over
$(C,X_{0:L}^{\mathrm{old}},J)$.  The next state
$Y=X_{J+1}^{\mathrm{old}}$ must be the actual transition sampled in this rollout.

Condition on one $(C,J,X_J^{\mathrm{old}})$.  The old, teacher, and current one-step kernels share
a positive, model-independent covariance:
\begin{equation}
p_j(y)=\N(y\mid\mu_j,\Sigma),
\qquad
j\in\{\mathrm{old},T,\theta\},
\qquad
\Sigma\succ0.
\label{eq:a3-kernels}
\end{equation}
Define
\begin{equation}
q_\alpha(y)
=(1-\alpha)p_{\mathrm{old}}(y)+\alpha p_T(y),
\qquad
Y\sim p_{\mathrm{old}},
\label{eq:a3-mixture}
\end{equation}
and
\begin{equation}
\gamma(Y)
=
\frac{\alpha p_T(Y)}{q_\alpha(Y)}
=
\sigma\!\left(
  \logit\alpha
  +\log\frac{p_T(Y)}{p_{\mathrm{old}}(Y)}
\right).
\label{eq:a3-responsibility}
\end{equation}
This is the flow-transition version of TOP-D's external probability-mixture
teacher~\cite{topd}.

\begin{proposition}[Exact local mixture-KL gradient]
\label{prop:a3-exact-gradient}
Let
$\mathcal K(\mu_\theta)=\KL(p_\theta\Vert q_\alpha)$.  At the behavior mean,
\begin{equation}
\left.
\nabla_{\mu_\theta}\mathcal K
\right|_{\mu_\theta=\mu_{\mathrm{old}}}
=
\E_{p_{\mathrm{old}}}[\gamma(Y)]
\Sigma^{-1}(\mu_{\mathrm{old}}-\mu_T).
\label{eq:a3-exact-gradient}
\end{equation}
\end{proposition}

\begin{proof}
Write $Y=\mu_\theta+\Sigma^{1/2}Z$ with $Z\sim\N(0,I)$.  Under this
reparameterization, $\log p_\theta(Y)$ is independent of $\mu_\theta$, so
\[
\nabla_{\mu_\theta}\mathcal K
=-\E_{p_\theta}[\nabla_y\log q_\alpha(Y)].
\]
The mixture score is
\[
\nabla_y\log q_\alpha(y)
=
\Sigma^{-1}\!\left(
  (1-\gamma(y))\mu_{\mathrm{old}}
  +\gamma(y)\mu_T-y
\right).
\]
At $\mu_\theta=\mu_{\mathrm{old}}$, the expectation of
$Y-\mu_{\mathrm{old}}$ vanishes, leaving \eqref{eq:a3-exact-gradient}.
\end{proof}

Now define the detached local surrogate
\begin{equation}
\mathcal S_{\mathrm{sum}}(\mu_\theta)
=
\E_{Y\sim p_{\mathrm{old}}}\!\left[
  \sg(\gamma(Y))
  \frac12
  \norm{\mu_\theta-\mu_T}_{\Sigma^{-1}}^2
\right].
\label{eq:a3-local-surrogate}
\end{equation}
Its gradient at $\mu_\theta=\mu_{\mathrm{old}}$ is exactly
\eqref{eq:a3-exact-gradient}.  If $\mu_\theta$ is a network output, both gradients are
left-multiplied by the same Jacobian $\partial_\theta\mu_\theta^\top$.

The complete Arm A3 event-sum loss averages this conditional surrogate over the actual behavior
SDE trajectory:
\begin{equation}
\mathcal L_{\mathrm{A3,sum}}(\theta)
=
\E_{\mathrm{SDE-old}}\!\left[
  \sg(\gamma_J(Y))
  \frac12
  \norm{
    \mu_{\theta,J}(X_J^{\mathrm{old}},C)
    -\mu_{T,J}(X_J^{\mathrm{old}},C)
  }_{\Sigma_J^{-1}}^2
\right].
\label{eq:a3-arm-loss}
\end{equation}
The fixed-state theorem applies inside this outer expectation.  Replacing
$\E_{\mathrm{SDE-old}}$ by the ODE occupation measure \eqref{eq:old-occupancy} would remove the
sampled stochastic transition and invalidate the density-ratio construction.

An event-mean implementation uses
\begin{equation}
\mathcal S_{\mathrm{mean}}=\frac{1}{D}\mathcal S_{\mathrm{sum}},
\qquad
\nabla\mathcal S_{\mathrm{mean}}
=
\frac1D\nabla\mathcal K
\quad\text{at }\mu_\theta=\mu_{\mathrm{old}}.
\label{eq:a3-mean-gradient}
\end{equation}
It matches the gradient of the dimension-normalized objective $\mathcal K/D$, not the numerical
gradient of the unnormalized joint KL.

\subsection{Ratio law and overlap}

\begin{proposition}[Gaussian ratio law]
\label{prop:a3-ratio-law}
Fix $\alpha\in(0,1)$ and let
\begin{equation}
\Delta=\mu_T-\mu_{\mathrm{old}},
\qquad
K=\frac12\Delta^\top\Sigma^{-1}\Delta
=\KL(p_{\mathrm{old}}\Vert p_T).
\label{eq:a3-K}
\end{equation}
Under $Y\sim p_{\mathrm{old}}$,
\begin{equation}
\log\frac{p_T(Y)}{p_{\mathrm{old}}(Y)}
\sim\N(-K,2K).
\label{eq:a3-ratio-law}
\end{equation}
Moreover,
\begin{equation}
\E_{p_{\mathrm{old}}}[\gamma(Y)]
\le
\min\!\left\{
\alpha,\,
\frac12\sqrt{\frac{\alpha}{1-\alpha}}e^{-K/4}
\right\}.
\label{eq:a3-overlap-bound}
\end{equation}
\end{proposition}

\begin{proof}
Write $Y=\mu_{\mathrm{old}}+\Sigma^{1/2}Z$.  Direct expansion gives
\[
\log\frac{p_T(Y)}{p_{\mathrm{old}}(Y)}
=Z^\top\Sigma^{-1/2}\Delta-K.
\]
The first term has variance $\Delta^\top\Sigma^{-1}\Delta=2K$.  For the second claim, put
$r=p_T/p_{\mathrm{old}}$.  The arithmetic--geometric mean inequality gives
\[
\gamma
=\frac{\alpha r}{1-\alpha+\alpha r}
\le
\frac12\sqrt{\frac{\alpha}{1-\alpha}}\sqrt r.
\]
Integrating under $p_{\mathrm{old}}$ gives the Gaussian Bhattacharyya coefficient
\[
\int\sqrt{p_{\mathrm{old}}p_T}
=
\exp\!\left(-\frac18\Delta^\top\Sigma^{-1}\Delta\right)
=e^{-K/4}.
\]
Finally, concavity of
$r\mapsto\alpha r/(1-\alpha+\alpha r)$ and $\E_{p_{\mathrm{old}}}[r]=1$
give the bound by $\alpha$.
\end{proof}

Equation~\eqref{eq:a3-overlap-bound} makes the mechanism explicit: the expected teacher pull is
controlled by overlap between the two complete transition events, not merely by per-coordinate
agreement.

\subsection{The ODE limit: the gate disappears, the control cost survives}

Keep $\alpha\in(0,1)$ fixed and let
\begin{equation}
\mu_j=x+hF_j,
\qquad
\Sigma_\varepsilon=\varepsilon hA,
\qquad
A\succ0 .
\label{eq:a3-euler-sde}
\end{equation}
Then
\begin{equation}
K_\varepsilon
=
\frac{h}{2\varepsilon}
\norm{F_T-F_{\mathrm{old}}}_{A^{-1}}^2.
\label{eq:a3-small-noise-K}
\end{equation}
If $F_T\neq F_{\mathrm{old}}$ and $h/\varepsilon\to\infty$, then
\begin{equation}
\gamma(Y)\to0
\quad\text{in probability},\qquad
\E[\gamma(Y)]
=
O\!\left(
  \exp\!\left[
    -\frac{h}{8\varepsilon}
    \norm{F_T-F_{\mathrm{old}}}_{A^{-1}}^2
  \right]
\right).
\label{eq:a3-small-noise-collapse}
\end{equation}
The mean-gradient with respect to the drift is
\begin{equation}
\left.
\nabla_{F_\theta}\mathcal K
\right|_{F_\theta=F_{\mathrm{old}}}
=
\frac{h}{\varepsilon}\E[\gamma(Y)]
A^{-1}(F_{\mathrm{old}}-F_T),
\label{eq:a3-drift-gradient}
\end{equation}
so the exponential overlap loss dominates the polynomial $1/\varepsilon$ factor and the teacher
gradient vanishes.

This limit requires $K_\varepsilon\to\infty$.  Taking $h\to0$ first at fixed nonzero diffusion
intensity instead makes the one-step $K_\varepsilon$ small.  The relevant ODE statement is the
zero-noise limit at a fixed practical solver grid, or any joint limit with $h/\varepsilon\to\infty$.

In contrast, the pairwise KL obeys
\begin{equation}
\frac{\varepsilon}{h}
\KL\!\left(
  \N(x+hF_\theta,\varepsilon hA)
  \,\middle\Vert\,
  \N(x+hF_T,\varepsilon hA)
\right)
=
\frac12\norm{F_\theta-F_T}_{A^{-1}}^2,
\label{eq:a3-pairwise-control}
\end{equation}
which is the control-energy geometry of \eqref{eq:control-energy-limit}.  The transition-mixture
KL itself instead approaches a flat component-selection cost:
\begin{equation}
\KL(p_{\mathrm{old}}\Vert q_\alpha)
\longrightarrow
-\log(1-\alpha).
\label{eq:a3-mixture-flat-limit}
\end{equation}
To verify the limit, put $r=p_T/p_{\mathrm{old}}$ and write
\begin{equation}
\KL(p_{\mathrm{old}}\Vert q_\alpha)
=
-\log(1-\alpha)
-\E_{p_{\mathrm{old}}}
\log\!\left(
  1+\frac{\alpha}{1-\alpha}r
\right).
\label{eq:a3-mixture-limit-expansion}
\end{equation}
Using $\log(1+x)\le2\sqrt{x}$ and the Bhattacharyya coefficient from
Proposition~\ref{prop:a3-ratio-law},
\begin{equation}
0
\le
-\log(1-\alpha)
-\KL(p_{\mathrm{old}}\Vert q_\alpha)
\le
2\sqrt{\frac{\alpha}{1-\alpha}}e^{-K_\varepsilon/4},
\label{eq:a3-mixture-limit-bound}
\end{equation}
which tends to zero when $K_\varepsilon\to\infty$.
Thus the quadratic metric has a useful deterministic limit, while the probability-mixture gate
does not.

\subsection{Temperature, latent sum, and latent mean}

For equal-covariance Gaussians and $T_g>0$,
\begin{equation}
\frac{1}{T_g}
\log\frac{\N(y\mid\mu_T,\Sigma)}
         {\N(y\mid\mu_{\mathrm{old}},\Sigma)}
=
\log\frac{\N(y\mid\mu_T,T_g\Sigma)}
         {\N(y\mid\mu_{\mathrm{old}},T_g\Sigma)}.
\label{eq:a3-temperature-covariance}
\end{equation}
Therefore dividing the log ratio by $T_g$ is pointwise equivalent to inflating both component
covariances by $T_g$.  A fully consistent inflated-covariance mixture would also sample
$Y\sim\N(\mu_{\mathrm{old}},T_g\Sigma)$ and divide the local quadratic by $T_g$.  Changing only
the gate temperature while retaining original samples and loss normalization is a tempered
surrogate, not the exact gradient of the original mixture KL.

Define the whitened per-coordinate gap
\begin{equation}
w^2
=
\frac1D
\Delta^\top\Sigma^{-1}\Delta .
\label{eq:a3-whitened-rms}
\end{equation}
Then
\begin{equation}
K=\frac{D}{2}w^2,
\qquad
L_{\mathrm{sum}}
\sim
\N\!\left(-\frac{D}{2}w^2,Dw^2\right),
\label{eq:a3-sum-law}
\end{equation}
while
\begin{equation}
L_{\mathrm{mean}}
=\frac{L_{\mathrm{sum}}}{D}
\sim
\N\!\left(-\frac12w^2,\frac{w^2}{D}\right).
\label{eq:a3-mean-law}
\end{equation}
The sum is the exact joint event likelihood.  The mean is exactly gate temperature $T_g=D$, or
covariance inflation by $D$.  As $D$ grows at fixed $w$, the sum gate separates exponentially;
the mean gate converges to
\begin{equation}
\sigma\!\left(\logit\alpha-\frac12w^2\right)
\label{eq:a3-mean-limit}
\end{equation}
and becomes deterministic across samples.  Intermediate temperatures define surrogates with an
explicit effective dimension.  This reproduces the companion exact-sum gate analysis.

\begin{tcolorbox}[colback=red!4,colframe=red!45!black,title=Exact boundary of Arm A3]
The local-gradient equality requires: the same complete latent event, positive shared
model-independent covariance, an actual transition sampled by the frozen old student,
$\mu_\theta=\mu_{\mathrm{old}}$ at the expansion point, detached old/teacher/gate quantities, and
the untempered joint ratio.  It is an equality of first derivatives at one on-policy point, not
an equality of loss values, Hessians, multi-step trajectory objectives, or gradients after
multiple stale-data updates.
\end{tcolorbox}

%% ============================================================
\section{Arm A4: Marginal-Mixture Conditional Flow Matching}
\label{sec:a4}
%% ============================================================

\subsection{Mixture marginals and their velocity}

Fix a conditioning value $C=c$.  All densities and fields in this section are conditional on
$c$, which is suppressed from the notation; the final training loss also averages over $C$.
Assume the old and teacher ODE marginals have densities satisfying their continuity equations:
\begin{equation}
\partial_t\rho_{b,t}
+\nabla_x\!\cdot(\rho_{b,t}v_b)=0,
\qquad
\rho_{b,0}=\rho_{\mathrm{base}},
\qquad
b\in\{\mathrm{old},T\}.
\label{eq:a4-source-continuity}
\end{equation}
Define
\begin{align}
q_t
&=(1-\alpha)\rho_{\mathrm{old},t}+\alpha\rho_{T,t},
\label{eq:a4-mixture-density}\\
q_tu_t
&=(1-\alpha)\rho_{\mathrm{old},t}\vold
  +\alpha\rho_{T,t}\vT .
\label{eq:a4-mixture-flux}
\end{align}

\begin{proposition}[Exact mixture continuity equation]
\label{prop:a4-continuity}
The pair $(q_t,u_t)$ satisfies
\begin{equation}
\partial_tq_t+\nabla_x\!\cdot(q_tu_t)=0,
\qquad
q_0=\rho_{\mathrm{base}}.
\label{eq:a4-continuity}
\end{equation}
Where $q_t(x)>0$, define
\begin{equation}
\gamma_t(x)
=
\frac{\alpha\rho_{T,t}(x)}{q_t(x)}.
\label{eq:a4-marginal-responsibility}
\end{equation}
Then
\begin{equation}
u_t(x)
=(1-\gamma_t(x))\vold(t,x)
+\gamma_t(x)\vT(t,x).
\label{eq:a4-mixture-velocity}
\end{equation}
\end{proposition}

\begin{proof}
Substitute \eqref{eq:a4-mixture-density}--\eqref{eq:a4-mixture-flux}:
\begin{align*}
\partial_tq_t+\nabla\!\cdot(q_tu_t)
&=(1-\alpha)
  [\partial_t\rho_{\mathrm{old},t}
   +\nabla\!\cdot(\rho_{\mathrm{old},t}\vold)]\\
&\quad
  +\alpha
  [\partial_t\rho_{T,t}
   +\nabla\!\cdot(\rho_{T,t}\vT)]
=0.
\end{align*}
The common base law gives $q_0=\rho_{\mathrm{base}}$.  Dividing the flux by $q_t$ gives
\eqref{eq:a4-mixture-velocity}.
\end{proof}

This is the deterministic analogue of a probability-mixture responsibility, but the
responsibility is over time marginals rather than one-step Dirac kernels.

\subsection{Two-source CFM without density ratios}

Sample one source branch for a whole trajectory and a training time:
\begin{equation}
B\sim\operatorname{Bernoulli}(\alpha),
\qquad
t\sim\omega,
\qquad
X_t=\Phi_t^B(Z),
\qquad
V_t=
\begin{cases}
\vold(t,X_t), & B=0,\\
\vT(t,X_t), & B=1.
\end{cases}
\label{eq:a4-branch-sampling}
\end{equation}
The practical objective is
\begin{equation}
\mathcal L_{\mathrm{A4}}(\theta)
=
\E_{C,t,B,Z}
\norm{\vtheta(t,X_t)-V_t}^2 .
\label{eq:a4-two-source-loss}
\end{equation}
The norm in \eqref{eq:a4-two-source-loss} is an event sum.  A per-dimension implementation divides
the entire loss, including the constant below, by $D$.

\begin{proposition}[Two-source CFM identity]
\label{prop:a4-cfm}
Under \eqref{eq:a4-branch-sampling},
$X_t\sim q_t$,
\begin{equation}
\Pr(B=1\mid X_t=x,t)=\gamma_t(x),
\qquad
\E[V_t\mid X_t=x,t]=u_t(x),
\label{eq:a4-conditional-mean}
\end{equation}
and
\begin{equation}
\mathcal L_{\mathrm{A4}}(\theta)
=
\E_{C,t,X_t\sim q_t(\cdot\mid C)}
\norm{\vtheta(t,X_t)-u_t(X_t)}^2
+\mathcal C_{\mathrm{var}},
\label{eq:a4-cfm-decomposition}
\end{equation}
where
\begin{equation}
\mathcal C_{\mathrm{var}}
=
\E_{C,t,X_t\sim q_t(\cdot\mid C)}\!\left[
  \gamma_t(X_t)(1-\gamma_t(X_t))
  \norm{\vT(t,X_t)-\vold(t,X_t)}^2
\right]
\label{eq:a4-variance-constant}
\end{equation}
is independent of $\theta$.
\end{proposition}

\begin{proof}
Marginalizing the branch label gives
$X_t\sim(1-\alpha)\rho_{\mathrm{old},t}+\alpha\rho_{T,t}=q_t$.  Bayes' rule gives
\eqref{eq:a4-marginal-responsibility}, and the conditional mean of the two possible velocity
labels is \eqref{eq:a4-mixture-velocity}.  Finally write
\[
\vtheta-V_t=(\vtheta-u_t)+(u_t-V_t).
\]
Conditioned on $(t,X_t)$, the cross term vanishes because
$\E[V_t-u_t\mid t,X_t]=0$.  The conditional variance of a two-point label is
$\gamma_t(1-\gamma_t)\norm{\vT-\vold}^2$, proving the result.
\end{proof}

The analytic optimum contains the density responsibility $\gamma_t$, but training never evaluates
it.  Regression on samples from the joint law $(B,X_t,V_t)$ learns the conditional mean
implicitly.  This is the standard conditional-flow-matching mechanism~\cite{flowmatching}.

\subsection{When the generated marginals are exactly the mixture}

Suppose $q_t$ and the flux in \eqref{eq:a4-mixture-flux} are continuously differentiable, $q_t>0$
on the relevant region, and $u_t$ is continuous in time, Lipschitz in state, and has at most
linear growth.  Then the ODE
\begin{equation}
\dot X_t=u_t(X_t),
\qquad
X_0\sim\rho_{\mathrm{base}},
\label{eq:a4-target-ode}
\end{equation}
has a unique flow, and its density is the unique classical solution of
\eqref{eq:a4-continuity}.  Therefore
\begin{equation}
(\Phi_t^u)_\#\rho_{\mathrm{base}}=q_t
=(1-\alpha)\rho_{\mathrm{old},t}+\alpha\rho_{T,t}.
\label{eq:a4-exact-marginals}
\end{equation}
Under population optimization and sufficient model capacity, the square-loss minimizer equals
$u_t$ almost everywhere under $\omega(t)q_t(x)\dd t\dd x$.  Arm A4 realizes the mixture path
exactly when the fitted field admits a continuous/Lipschitz representative that equals $u_t$ on
the region traversed by the target flow; equality of abstract $L^2$ equivalence classes alone is
not enough to define an ODE.

These regularity conditions are not automatic.  Where both component densities are positive,
differentiating \eqref{eq:a4-marginal-responsibility} gives
\begin{equation}
\nabla\gamma_t
=
\gamma_t(1-\gamma_t)
\left(
  \nabla\log\rho_{T,t}
  -\nabla\log\rho_{\mathrm{old},t}
\right).
\label{eq:a4-gamma-gradient}
\end{equation}
Low-overlap distributions can create a thin, stiff boundary where $u_t$ switches between the two
source fields.

\subsection{Exact contraction statements}

\begin{proposition}[Contraction of the ideal marginal mixture]
\label{prop:a4-contraction}
Assume the displayed KL terms are finite and both source laws have finite $p$-th moments.  At
every time $t$,
\begin{align}
\norm{q_t-\rho_{T,t}}_1
&=(1-\alpha)
  \norm{\rho_{\mathrm{old},t}-\rho_{T,t}}_1,
\label{eq:a4-l1-contraction}\\
\KL(q_t\Vert\rho_{T,t})
&\le
(1-\alpha)
\KL(\rho_{\mathrm{old},t}\Vert\rho_{T,t}),
\label{eq:a4-kl-contraction}\\
W_p^p(q_t,\rho_{T,t})
&\le
(1-\alpha)
W_p^p(\rho_{\mathrm{old},t},\rho_{T,t}),
\qquad p\ge1.
\label{eq:a4-wasserstein-contraction}
\end{align}
The first line is equivalently an exact total-variation contraction.  For $p=1$,
\eqref{eq:a4-wasserstein-contraction} is an equality; for $p>1$ it need not be.
\end{proposition}

\begin{proof}
The signed density difference is
$q_t-\rho_{T,t}=(1-\alpha)(\rho_{\mathrm{old},t}-\rho_{T,t})$, proving
\eqref{eq:a4-l1-contraction}.  Convexity of $r\mapsto r\log r$ gives
\[
\left(
  (1-\alpha)\frac{\rho_{\mathrm{old},t}}{\rho_{T,t}}+\alpha
\right)
\log\!\left(
  (1-\alpha)\frac{\rho_{\mathrm{old},t}}{\rho_{T,t}}+\alpha
\right)
\le
(1-\alpha)
\frac{\rho_{\mathrm{old},t}}{\rho_{T,t}}
\log\frac{\rho_{\mathrm{old},t}}{\rho_{T,t}},
\]
which integrates against $\rho_{T,t}$ to \eqref{eq:a4-kl-contraction}.  For Wasserstein distance,
mix an optimal
coupling between $\rho_{\mathrm{old},t}$ and $\rho_{T,t}$ with weight $1-\alpha$ and the identity
coupling of $\rho_{T,t}$ with itself with weight $\alpha$.  Its transport cost is
$(1-\alpha)W_p^p(\rho_{\mathrm{old},t},\rho_{T,t})$, proving the upper bound.  The
Kantorovich--Rubinstein dual formula gives equality at $p=1$.
\end{proof}

These are properties of the ideal mixture $q_t$.  A finite-capacity student, if a global
population optimum is attained, returns the best $L^2(q_t)$ fit to $u_t$ and need not preserve
them exactly.  If the fitted field $\widehat v$ is $L$-Lipschitz, the same coupling argument as
Proposition~\ref{prop:gronwall} gives the practical error bound
\begin{equation}
W_2^2(\widehat\rho_t,q_t)
\le
t\int_0^t
e^{2L(t-s)}
\E_{X_s\sim q_s}
\norm{\widehat v(s,X_s)-u_s(X_s)}^2\dd s .
\label{eq:a4-capacity-bound}
\end{equation}

\subsection{Cost and the branch-selection distinction}

Exact A4 sampling requires states from both $\rho_{\mathrm{old},t}$ and $\rho_{T,t}$.  In
practice this means complete teacher trajectories for an $\alpha$ fraction of examples, or cached
teacher trajectories reused across sampled training times.  Evaluating the teacher field only on
old states does not create samples from the teacher marginal and reduces A4 back to an
old-occupancy field objective.

Choosing $B$ once per trajectory gives the path-law mixture used in
\eqref{eq:a4-branch-sampling}.  Choosing independent branches at every small Euler step instead
has conditional mean field
\begin{equation}
\bar v=(1-\alpha)\vold+\alpha\vT,
\label{eq:a4-rapid-switch}
\end{equation}
and converges, as switching variance vanishes with the step size, to the arithmetic-average ODE
of Arm A1.  In general,
\begin{equation}
\bar v-u_t
=(\alpha-\gamma_t)(\vT-\vold),
\label{eq:a4-average-gap}
\end{equation}
so the two constructions coincide only in special cases such as equal marginals or equal fields.

%% ============================================================
\section*{Synthesis}
\addcontentsline{toc}{section}{Synthesis}
%% ============================================================

\begin{center}
\begin{tabular}{@{}
  >{\raggedright\arraybackslash}p{0.08\textwidth}
  >{\raggedright\arraybackslash}p{0.23\textwidth}
  >{\raggedright\arraybackslash}p{0.27\textwidth}
  >{\raggedright\arraybackslash}p{0.31\textwidth}
@{}}
\toprule
Arm & mathematical object & exact statement & principal boundary\\
\midrule
A0 & old-occupancy field MSE
& exact weighted regression / finite-capacity least-squares fit
& no explicit trust region; no distribution objective\\
\addlinespace
A1 & quadratic field geometry
& exact field barycenter and trust-ball projection; local Euler transport view
& one on-policy step only scales A0; not a density mixture\\
\addlinespace
A2 & forward-target predictive risk
& exact posterior only for a declared shared-covariance residual model
& not a state-density ratio; calibration changes the surrogate\\
\addlinespace
A3 & stochastic one-step transition mixture
& exact mixture-KL first derivative at the old policy
& requires SDE covariance and $T_g=1$; high-dimensional overlap collapse\\
\addlinespace
A4 & mixture of ODE time marginals
& exact marginal path under population fitting and ODE regularity
& requires teacher trajectories, capacity, and a regular mixture velocity\\
\bottomrule
\end{tabular}
\end{center}

The five arms answer different questions rather than forming a single ranking.  A0 is the cheapest
ODE baseline.  A1 is the appropriate deterministic trust-region baseline.  A2 adds inexpensive,
solver-agnostic evidence about teacher predictive competence.  A3 is the strict transition-level
probability-mixture construction when an SDE is acceptable.  A4 is the strict deterministic
probability-mixture comparator at the marginal level, paid for with teacher trajectories.

For practical ODE distillation, the main comparison should therefore keep these roles separate:
\begin{enumerate}[leftmargin=1.6em]
\item compare A0 and one-step A1 at a matched effective teacher-gradient scale, because
      Proposition~\ref{prop:a1-first-step} predicts equality otherwise;
\item test A1 with an actually enforced field RMS budget or multiple frozen-old inner updates;
\item compare A2 against an ungated target with the same timestep and norm normalization;
\item retain A3 as the stochastic local-mixture reference and log its joint/per-dimension gap and
      responsibility saturation;
\item use A4 on a smaller scale as the exact marginal-mixture comparator.
\end{enumerate}

\begin{tcolorbox}[colback=gray!5,colframe=gray!55,title=Main formal conclusion]
There is no nontrivial deterministic continuation of the SDE transition-responsibility gate.
What transfers to ODEs is the rescaled quadratic control geometry (A1), not the likelihood ratio.
An exact deterministic probability mixture is still possible, but it lives at the level of ODE
time marginals and is trained by two-source conditional flow matching (A4).
\end{tcolorbox}

\begin{thebibliography}{9}

\bibitem{flowmatching}
Yaron Lipman, Ricky T. Q. Chen, Heli Ben-Hamu, Maximilian Nickel, and Matt Le.
\newblock Flow Matching for Generative Modeling.
\newblock \emph{International Conference on Learning Representations}, 2023.
\newblock \url{https://arxiv.org/abs/2210.02747}.

\bibitem{diffusionopd}
Quanhao Li, Junqiu Yu, Kaixun Jiang, et al.
\newblock DiffusionOPD: A Unified Perspective of On-Policy Distillation in Diffusion Models.
\newblock 2026.
\newblock \url{https://arxiv.org/abs/2605.15055}.

\bibitem{topd}
Zhengpeng Xie, Li Lyna Zhang, Zeke Xie, and Mao Yang.
\newblock Trust Region Policy Distillation.
\newblock 2026.
\newblock \url{https://arxiv.org/abs/2607.04751}.

\bibitem{flowdppo}
Bowen Ping, Xiangxin Zhou, Penghui Qi, Minnan Luo, Liefeng Bo, and Tianyu Pang.
\newblock Flow-DPPO: Divergence Proximal Policy Optimization for Flow Matching Models.
\newblock 2026.
\newblock \url{https://arxiv.org/abs/2606.11025}.

\end{thebibliography}

\end{document}
